% Part: computability
% Chapter: computability-theory
% Section: introduction

\documentclass[../../../include/open-logic-section]{subfiles}

\begin{document}

\olfileid{cmp}{thy}{int}
\olsection{ভূমিকা}

যুক্তিবিদ্যার যে শাখাটি \emph{গণনসাধ্যতা তত্ত্ব} নামে পরিচিত, সেখানে
অপেক্ষক ও সেটের গণনসাধ্যতা বা আপেক্ষিক গণনসাধ্যতা-সংক্রান্ত বিষয় নিয়ে
আলোচনা করা হয়। ক্লিনির প্রভাবের প্রমাণ হলো, বিষয়টি আগে
\emph{পুনরাবৃত্তি তত্ত্ব} নামে পরিচিত ছিল; আজকাল দুটি নামই প্রচলিত।

কোনো অপেক্ষক~$f\colon \Nat \pto \Nat$-কে \emph{আংশিক গণনসাধ্য} বলব,
যদি কোনো গণনা-মডেলে সেটিকে গণনা করা যায়। $f$ সর্বত্র-সংজ্ঞায়িত হলে
আমরা শুধু বলব, $f$~\emph{গণনসাধ্য}। যে সম্বন্ধ~$R$-এর চরিত্রসূচক
অপেক্ষক~$\Char{R}$ গণনসাধ্য, তাকেও গণনসাধ্য বলা হয়। $f$ ও~$g$ আংশিক
অপেক্ষক হলে, $f(x) \fdefined$ দিয়ে বোঝাব যে~$x$-এ $f$ সংজ্ঞায়িত,
অর্থাৎ~$x$ হলো~$f$-এর সংজ্ঞাক্ষেত্রের অন্তর্ভুক্ত; আর $f(x)
\fundefined$ দিয়ে বিপরীতটি, অর্থাৎ~$x$-এ $f$ অসংজ্ঞায়িত, বোঝাব।
$f(x) \simeq g(x)$ দিয়ে বোঝাব যে হয় $f(x)$ ও $g(x)$ উভয়ই
অসংজ্ঞায়িত, নয়তো উভয়ই সংজ্ঞায়িত ও সমান।

কোনো নির্দিষ্ট গণনা-মডেলের উল্লেখ না করেও গণনসাধ্যতার তত্ত্ব অনুসন্ধান
করা যায়। এজন্য দেখানো হয় যে একটি সার্বজনীন আংশিক গণনসাধ্য
অপেক্ষক~$\fn{Un}(k, x)$ আছে। এর সাহায্যে আংশিক গণনসাধ্য অপেক্ষকগুলিকে
তালিকায়িত করা যায়। $k$-তম একস্থানী আংশিক গণনসাধ্য অপেক্ষক বোঝাতে আমরা
সংকেত~$\cfind{k}$ গ্রহণ করব; এর সংজ্ঞা হলো
$\cfind{k}(x) \simeq \fn{Un}(k, x)$। (ক্লিনি এই কাজে $\{ k \}$ ব্যবহার
করেছিলেন, কিন্তু সম্প্রতি এই সংকেতের ব্যবহার কমেছে।) একটু সাধারণভাবে,
যেকোনো স্থানসংখ্যার আংশিক গণনসাধ্য অপেক্ষককে একই নিয়মে তালিকায়িত করা
যায়; $k$-তম $n$-স্থানী আংশিক পুনরাবৃত্ত অপেক্ষক বোঝাতে আমরা
$\cfind{k}[n]$ ব্যবহার করব।

$f(\vec x, y)$ সর্বত্র-সংজ্ঞায়িত বা আংশিক অপেক্ষক হলে,
$\umin{y}{f (\vec x, y)}$ হলো~$\vec x$-এর সেই অপেক্ষক, যা এমন ক্ষুদ্রতম~$y$
ফেরত দেয় যার জন্য $f(\vec x, y) = 0$, তবে ধরে নেওয়া হয় যে
$f(\vec x, 0)$, \dots, $f(\vec x, y-1)$ সবগুলি সংজ্ঞায়িত; এমন কোনো~$y$
না থাকলে $\umin{y}{f (\vec x, y)}$ অসংজ্ঞায়িত। $R(\vec x, y)$ কোনো
সম্বন্ধ হলে, $\umin{y}{R(\vec x, y)}$-কে এমন ক্ষুদ্রতম~$y$ হিসেবে
সংজ্ঞায়িত করা হয় যার জন্য $R(\vec x, y)$ সত্য; অন্যভাবে বললে, এমন
ক্ষুদ্রতম~$y$ যার জন্য $1 \tsub \Char{R}(\vec x, y) = 0$।

কোনো অপেক্ষক গণনসাধ্য দেখানোর দুটি উপায় আছে:
\begin{enumerate}
\item বিধিবদ্ধভাবে: একটি টুরিং যন্ত্র বা আংশিক পুনরাবৃত্ত অপেক্ষক
  প্রত্যক্ষভাবে বর্ণনা করে দেখানো যে সেটি অভীষ্ট অপেক্ষকটিকেই গণনা করে;
\item অনানুষ্ঠানিকভাবে: সেটি গণনা করে এমন একটি অ্যালগরিদম বর্ণনা করে
  চার্চের থিসিস প্রয়োগ করা।
\end{enumerate}
দুই পদ্ধতির মধ্যে কোনো তীক্ষ্ণ সীমারেখা নেই; একটি অ্যালগরিদমের বিশদ
বর্ণনায় এতটুকু তথ্য থাকা উচিত, যাতে নীতিগতভাবে উপযুক্ত টুরিং যন্ত্র বা
আংশিক পুনরাবৃত্ত সংজ্ঞার অনুক্রম কীভাবে তৈরি করা যায় তা মোটামুটি স্পষ্ট
হয়। সম্পূর্ণ বিধিবদ্ধ সংজ্ঞা খুব তথ্যবহুল হওয়ার সম্ভাবনা কম; তাই আমরা
দুটি পদ্ধতির মধ্যে উপযুক্ত মধ্যপন্থা খুঁজব—অর্থাৎ এমন স্বজ্ঞাসম্মত অথচ
বিধিবদ্ধ প্রমাণ দেব, যাতে নির্ভুল সংজ্ঞাগুলি যে পাওয়া যেতে পারে তা
প্রতিষ্ঠিত হয়।

\end{document}
